\documentclass[11pt, a4paper]{article}

% --- UNIVERSAL PREAMBLE BLOCK ---
\usepackage[a4paper, top=2.5cm, bottom=2.5cm, left=2cm, right=2cm]{geometry}
\usepackage{fontspec} % Required for Noto Sans/Babel
\usepackage{amsmath} % For equation formatting
\usepackage{tabularx} % For X column type
\usepackage{booktabs} % For professional table lines (\toprule, \midrule, \bottomrule)

\usepackage[english, bidi=basic, provide=*]{babel}

\babelprovide[import, onchar=ids fonts]{english}

% Set default/Latin font to Sans Serif
\babelfont{rm}{Noto Sans}

% --- END UNIVERSAL PREAMBLE BLOCK ---

\title{ABL Closure Project Glossary}
\author{Machine Learning Assisted SBL Closure Team}
\date{\today}

\begin{document}

\begin{table}[t]
\centering
\renewcommand{\arraystretch}{1.35} % Increased stretch for better readability with formulas
\caption{Canonical definitions and symbols used in the ABL Closure Project and associated machine-learning components.}
\label{tab:abl-glossary}
\begin{tabularx}{\linewidth}{l c X}
\toprule
\textbf{Term} & \textbf{Symbol} & \textbf{Definition / Project Context} \\
\midrule

Obukhov length & $L$ & Scale at which shear and buoyancy production balance; fundamental length scale for SBL stability. \\
\addlinespace

Dimensionless height & $\zeta$ & $\zeta = z/L$. Non-dimensional height used in MOST functions; $\zeta>0$ in stable conditions. \\
\addlinespace

Geometric mean height & $z_g$ & $z_g = \sqrt{z_0 z_1}$. Representative height for log-layer profiles and where $\mathrm{Ri}_g$ is often evaluated. \\
\addlinespace

Logarithmic mean height & $z_L$ & $z_L = (z_1 - z_0) / \ln(z_1 / z_0)$. Exact representative height for the velocity difference in a logarithmic wind law. \\
\addlinespace

Gradient Richardson number & $\mathrm{Ri}_g$ & $\displaystyle \mathrm{Ri}_g(z)=\frac{(g/\theta)\,(\partial\theta/\partial z)}{(\partial U/\partial z)^2}$. Local stability diagnostic used to determine the collapse threshold $\mathrm{Ri}_c$. \\
\addlinespace

Bulk Richardson number & $\mathrm{Ri}_b$ & $\displaystyle \mathrm{Ri}_b=\frac{(g/\bar{\theta})\,(\Delta\theta\,\Delta z)}{(\Delta U)^2}$. Layer-averaged stability measure and main input for NWP closures. \\
\addlinespace

Bias ratio & $B$ & $B = \mathrm{Ri}_g(z_g) / \mathrm{Ri}_b$. Diagnostic of curvature-induced bias. $B>1$ indicates that $\mathrm{Ri}_b$ underestimates local stability. \\
\addlinespace

Grid damping factor & $G$ & $G(\zeta,\Delta z)$; multiplicative correction applied to eddy diffusivity: $K_{\text{new}} = K_{\text{old}} \cdot G$. $\text{ML}$ model predicts $G$ from $(\zeta,\Delta z)$. \\
\addlinespace

Eddy diffusivities & $K_m, K_h$ & Momentum and heat diffusivities ($m^2 s^{-1}$). Grid damping factor $G$ is applied directly. \\
\addlinespace

Critical Richardson number & $\mathrm{Ri}_c,\ \mathrm{Ri}_c^\ast$ & Threshold for turbulence collapse. $\mathrm{Ri}_c^\ast$ denotes dynamic, state-dependent value (the primary target for symbolic regression). \\
\addlinespace

Neutral curvature invariant & $\Delta$ & $\Delta = a_h - 2 a_m$. Determines the sign of $\partial_\zeta^2 \mathrm{Ri}_g$ near neutrality, which sets the initial sign of the bias ratio $B$. \\
\addlinespace

Turbulent kinetic energy & $\text{TKE}$ & $0.5(\overline{u'^2}+\overline{v'^2}+\overline{w'^2})$. Used in higher-order closures for turbulence memory and intermittency. \\

\bottomrule
\end{tabularx}
\end{table}

\end{document}